\documentclass[11pt,a4paper]{article}

\usepackage[utf8]{inputenc}
\usepackage[T1]{fontenc}
\usepackage[french]{babel}
\usepackage{lmodern}
\usepackage{microtype}
\usepackage[margin=2.3cm]{geometry}
\usepackage{booktabs}
\usepackage{tabularx}
\usepackage{array}
\usepackage{xcolor}
\usepackage{enumitem}
\usepackage{float}
\usepackage{amsmath}
\usepackage[round,authoryear]{natbib}
\usepackage[colorlinks=true,linkcolor=teal!70!black,citecolor=teal!70!black,urlcolor=teal!70!black]{hyperref}

\newcolumntype{Y}{>{\raggedright\arraybackslash}X}
\setlist{itemsep=2pt,topsep=4pt}

\title{\textbf{Livrable L2 -- Cahier expérimental}\\[4pt]
\large Allocation adaptative des ressources de calcul et de recherche documentaire\\ pour des systèmes LLM locaux, économes et fiables\\[2pt]
\normalsize\textit{Benchmark, systèmes comparés, métriques, critères de succès et pilote}}
\author{Alae Lahbichi -- Étudiant en 5\textsuperscript{e} année IA/Data\\[2pt]
\small\href{mailto:a.lahbichi2004@gmail.com}{\texttt{a.lahbichi2004@gmail.com}}\\
\small Dépôt GitHub : \href{https://github.com/AlaeLahbichi/Adaptive-Compute-and-RAG-for-Efficient-Reliable-Local-LLMs}{\footnotesize\texttt{AlaeLahbichi/Adaptive-Compute-and-RAG-for-Efficient-Reliable-Local-LLMs}}}
\date{Version 1.0-pilote -- 18 octobre 2026}

\begin{document}
\maketitle

\begin{center}
\fcolorbox{teal!60!black}{teal!5}{\parbox{0.92\textwidth}{\small
\textbf{Rôle du document.} Ce cahier fixe, \emph{avant} toute expérience finale, le benchmark, les
systèmes comparés, les métriques, la méthode statistique et les critères de succès. Les paramètres
sont consignés dans \texttt{configs/protocole\_v1.json}. Après validation du pilote
(section~\ref{sec:pilote}), le protocole sera gelé par l'étiquette Git \texttt{protocole-v1.0} :
toute modification ultérieure devra être datée, justifiée et signalée dans l'article.}}
\end{center}

\begin{abstract}
\noindent Ce cahier expérimental définit le cadre d'évaluation d'une politique qui alloue, requête par
requête, la taille du modèle, la recherche documentaire et l'abstention. Le benchmark réunit
630 questions en trois catégories (simple, complexe, non couverte), tirées de SQuAD~2.0 et HotpotQA,
et un corpus de 3\,293 passages. Il est séparé en un jeu de développement (180 questions) et un jeu de
test (450 questions) réservé à l'évaluation finale. Deux modèles déjà présents sur la machine sont
utilisés : Qwen2.5-1.5B et Qwen3-4B, quantifiés en 4 bits et exécutés avec MLX sur un Apple M3 Pro.
Quatre stratégies fixes servent de référence. Le critère de succès principal est une réduction d'au
moins 20\,\% de la latence moyenne par rapport à la stratégie « grand modèle + RAG », pour une perte
de score global d'au plus 2 points. Un pilote sur le jeu de développement valide la chaîne complète
et estime le temps de calcul.
\end{abstract}

\clearpage
\tableofcontents
\clearpage

%======================================================================
\section{Objet et principes}\label{sec:objet}
%======================================================================

Le cahier des charges impose de \emph{définir le succès avant les tests} et de \emph{séparer
développement et test}. Trois principes en découlent :
\begin{enumerate}
\item \textbf{Pré-enregistrement.} Les questions, les métriques, les systèmes de référence et les seuils
de succès sont fixés ici avant de développer la méthode adaptative (Partie IV).
\item \textbf{Étanchéité du test.} Les seuils de la méthode sont réglés uniquement sur le jeu de
développement. Le jeu de test n'est exécuté qu'une fois, sur la version finale, après le gel.
\item \textbf{Mêmes conditions pour tous.} Tous les systèmes sont évalués sur les mêmes questions, sur la
même machine, avec les mêmes modèles, prompts et paramètres de génération.
\end{enumerate}

%======================================================================
\section{Environnement matériel et logiciel}\label{sec:env}
%======================================================================

\begin{table}[H]
\centering\small
\begin{tabularx}{\textwidth}{l Y}
\toprule
\textbf{Élément} & \textbf{Valeur} \\
\midrule
Machine & Apple M3 Pro, 12 cœurs, 18 Go de mémoire unifiée, macOS 27.0.1 \\
Exécution des modèles & MLX 0.32.3 et mlx-lm 0.32.0 (Python 3.12.11) \\
Petit modèle & \texttt{mlx-community/Qwen2.5-1.5B-Instruct-4bit} (environ 1,0 Go en mémoire) \\
Grand modèle & \texttt{mlx-community/Qwen3-4B-Instruct-2507-4bit} (environ 2,5 Go en mémoire) \\
Embeddings & \texttt{sentence-transformers/paraphrase-multilingual-MiniLM-L12-v2} \\
Recherche & Similarité cosinus exacte en NumPy, équivalente à un index FAISS \texttt{IndexFlatIP} \\
Énergie & \texttt{powermetrics} (macOS, droits administrateur), échantillonnage toutes les 200~ms \\
\bottomrule
\end{tabularx}
\caption{Environnement d'exécution. Les révisions exactes des modèles et les versions des paquets sont
relevées automatiquement par \texttt{bench/check\_env.py} à chaque série d'expériences.}
\label{tab:env}
\end{table}

\paragraph{Choix de MLX plutôt qu'Ollama.} Le cahier des charges cite Ollama parmi les options
possibles. MLX est retenu parce qu'il est optimisé pour Apple Silicon, que les modèles étaient déjà
présents localement (aucun téléchargement supplémentaire) et qu'il donne accès au pic de mémoire
consommé par le modèle. Les deux modèles tiennent ensemble en mémoire (environ 3,5 Go), ce qui permet
de passer de l'un à l'autre sans rechargement.

\paragraph{Choix des modèles.} Le ratio de taille est d'environ 2,7 (1,5 contre 4 milliards de
paramètres). Les deux modèles appartiennent à deux générations de Qwen : l'écart de qualité ne
s'explique donc pas uniquement par la taille. Ce point est sans conséquence pour l'objectif, qui est
de disposer de deux niveaux de coût et de capacité, mais il sera mentionné comme limite. Llama-3.2-3B,
également disponible, pourra servir de test de robustesse (QR4).

%======================================================================
\section{Benchmark}\label{sec:benchmark}
%======================================================================

\subsection{Sources}
\begin{itemize}
\item \textbf{SQuAD 2.0}, jeu de développement \citep{rajpurkar2018squad2} : 35 articles de Wikipédia
(histoire, sciences, droit, géographie, institutions), avec des questions à réponse et des questions
sans réponse rédigées pour ressembler aux premières.
\item \textbf{HotpotQA}, jeu de développement en configuration \emph{distractor} \citep{yang2018hotpotqa} :
questions qui exigent de combiner deux pages Wikipédia, accompagnées de 8 paragraphes distracteurs.
\end{itemize}
Les deux jeux sont sous licence CC BY-SA 4.0. Leur empreinte SHA-256 est vérifiée à la construction.

\subsection{Corpus documentaire}
Le corpus est la « bibliothèque » dans laquelle la recherche documentaire puise ; il ne sert pas à
entraîner les modèles. Il réunit les 1\,204 paragraphes de SQuAD 2.0 et les 2\,089 paragraphes
distincts associés aux questions HotpotQA retenues, soit \textbf{3\,293 passages}. Toutes les
questions interrogent ce même corpus : la recherche doit donc retrouver les bons passages parmi
plusieurs milliers.

\subsection{Catégories de questions}
\begin{table}[H]
\centering\small
\begin{tabularx}{\textwidth}{l l Y l}
\toprule
\textbf{Catégorie} & \textbf{Source} & \textbf{Définition} & \textbf{Attendu} \\
\midrule
Simple & SQuAD 2.0 & La réponse se trouve dans un seul passage & Répondre \\
Complexe & HotpotQA & La réponse exige de combiner deux passages & Répondre \\
Non couverte & SQuAD 2.0 & Le passage d'origine ne contient pas la réponse & S'abstenir \\
Ambiguë & Annotation manuelle & Plusieurs interprétations possibles & Préciser ou couvrir \\
\bottomrule
\end{tabularx}
\caption{Catégories du benchmark.}
\label{tab:categories}
\end{table}

\paragraph{Catégorie « ambiguë ».} Aucun des deux jeux ne fournit de questions ambiguës alignées sur le
corpus. Elles seront rédigées à la main, avec un objectif de 50 questions, dans le fichier
\texttt{data/benchmark/ambigues\_a\_annoter.jsonl}, qui contient un exemple et le format attendu. Elles
seront évaluées séparément et ne font pas partie du critère de succès principal.

\subsection{Taille et séparation}
\begin{table}[H]
\centering\small
\begin{tabular}{l c c c c}
\toprule
& \textbf{Simple} & \textbf{Complexe} & \textbf{Non couverte} & \textbf{Total} \\
\midrule
Développement & 60 & 60 & 60 & 180 \\
Test & 150 & 150 & 150 & 450 \\
\bottomrule
\end{tabular}
\caption{Répartition des questions (graine aléatoire 2026).}
\label{tab:tailles}
\end{table}

Les questions SQuAD sont séparées \emph{par article} (30\,\% des articles pour le développement) afin
qu'aucun paragraphe ne serve à la fois au réglage et au test. Les questions HotpotQA sont tirées sans
recouvrement. Avec 150 questions par catégorie, l'intervalle de confiance à 95\,\% d'une proportion
proche de 50\,\% a une demi-largeur d'environ 8 points par catégorie, et d'environ 4,6 points sur
l'ensemble du test.

\subsection{Format d'une question}
Chaque question est une ligne JSON :
\begin{small}
\begin{verbatim}
{"id": "hotpot:5a8b57f2...", "split": "dev", "categorie": "complexe",
 "source": "hotpotqa", "question": "...", "reponses": ["..."],
 "passages_gold": ["hotpot:3f2a...", "hotpot:9c1d..."],
 "comportement_attendu": "repondre", "type_hotpot": "bridge"}
\end{verbatim}
\end{small}

\subsection{Limites connues du benchmark}
\begin{itemize}
\item Certaines questions de SQuAD sont formulées pour un paragraphe précis et deviennent floues dans un
grand corpus (par exemple « What was the source of the mistake? »). Elles pénalisent toutes les
stratégies de la même manière, mais abaissent les scores absolus.
\item Une question « non couverte » l'est par rapport à son paragraphe d'origine ; il est possible,
quoique rare, qu'un autre passage ou la connaissance interne du modèle contienne une réponse.
\item Certaines réponses de référence de HotpotQA sont bruitées (par exemple « Mary Harron (born January
12, 1953) ») ; le F1 atténue ce problème sans le supprimer.
\item Les données sont en anglais et relèvent de la culture générale. La généralisation à un corpus
spécialisé (QR4) devra être testée séparément.
\end{itemize}

%======================================================================
\section{Systèmes comparés}\label{sec:systemes}
%======================================================================

\subsection{Stratégies fixes (références)}
\begin{table}[H]
\centering\small
\begin{tabularx}{\textwidth}{l l c Y}
\toprule
\textbf{Code} & \textbf{Modèle} & \textbf{Passages} & \textbf{Rôle} \\
\midrule
A & Grand & 0 & Baseline A du cahier des charges : le modèle puissant répond seul \\
B & Grand & 3 & \textbf{Référence principale} : RAG systématique avec le grand modèle \\
C & Petit & 0 & Baseline C du cahier des charges : le petit modèle répond seul \\
D & Petit & 3 & RAG systématique avec le petit modèle \\
\bottomrule
\end{tabularx}
\caption{Stratégies fixes. B correspond à la configuration « tout au maximum » qu'un système non
adaptatif utiliserait par défaut ; c'est la référence des critères de succès.}
\label{tab:strategies}
\end{table}

\subsection{Références issues de la littérature (Partie III)}
Conformément aux conclusions du livrable L1, deux références supplémentaires seront ajoutées :
\begin{itemize}
\item un \textbf{routeur fondé sur la seule question} (classifieur léger, par exemple TF-IDF + régression
logistique), qui sert de contrôle pour vérifier que les signaux plus coûteux apportent réellement
quelque chose ;
\item une \textbf{réimplémentation simplifiée d'Adaptive-RAG}, avec un classifieur de complexité qui
choisit entre 0 passage, 3 passages, ou une recherche complémentaire de 6 passages.
\end{itemize}

\subsection{Espace de décision de la méthode adaptative (Partie IV)}
Pour chaque question, la politique choisit une action parmi :
\{petit, grand\} $\times$ \{aucune recherche, $k=3$, recherche complémentaire $k=6$\} $\times$
\{répondre, s'abstenir\}. Le coût de calcul de ses propres signaux (embeddings, scores, éventuelle
génération préalable) est inclus dans la latence mesurée.

%======================================================================
\section{Protocole d'exécution}\label{sec:execution}
%======================================================================

\begin{itemize}
\item \textbf{Prompts figés} (annexe~\ref{ann:prompts}), identiques pour tous les systèmes, avec une
consigne d'abstention explicite : répondre exactement \texttt{UNANSWERABLE} si la réponse est inconnue.
\item \textbf{Génération} gloutonne (température 0), au plus 32 tokens, avec le gabarit de conversation
de chaque modèle. Les sorties sont donc déterministes.
\item \textbf{Modèles préchargés} en mémoire. Trois requêtes de chauffe sont exécutées et ignorées avant
chaque stratégie.
\item \textbf{Ordre des questions} mélangé avec une graine fixe, identique pour toutes les stratégies.
\item \textbf{Latence} mesurée de bout en bout avec une horloge monotone : calcul des signaux, recherche
documentaire et génération.
\item \textbf{Journalisation} de chaque réponse en JSONL : sortie brute, passages utilisés, latence,
tokens, note.
\item \textbf{Répétitions} : les sorties étant déterministes, la qualité est mesurée une fois ; la
latence de l'évaluation finale est mesurée sur 3 passages complets, et c'est la médiane par question
qui est retenue.
\end{itemize}

%======================================================================
\section{Métriques}\label{sec:metriques}
%======================================================================

\begin{table}[H]
\centering\small
\begin{tabularx}{\textwidth}{l l Y}
\toprule
\textbf{Dimension} & \textbf{Métrique} & \textbf{Définition} \\
\midrule
Qualité & F1, exact match & Normalisation officielle de SQuAD ; maximum sur les réponses de référence \\
Qualité & \textbf{Score global} & Moyenne sur toutes les questions : F1 si une réponse est attendue ; 1 en cas d'abstention et 0 sinon pour une question non couverte \\
Abstention & Couverture & Part des questions auxquelles le système répond \\
Abstention & Risque & Part d'erreurs parmi les réponses émises (correcte si F1 $\geq$ 0,5) \\
Recherche & Rappel@$k$ & 1 si tous les passages justificatifs sont parmi les $k$ passages récupérés \\
Efficacité & Latence & Moyenne, médiane et 95\textsuperscript{e} centile (P95), en secondes \\
Efficacité & Tokens & Tokens du prompt et tokens générés, par question \\
Efficacité & Mémoire & Pic de mémoire MLX par stratégie \\
Énergie & Joules / question & Puissance CPU + GPU mesurée par \texttt{powermetrics}, intégrée sur la durée de la série, après soustraction de la consommation au repos \\
\bottomrule
\end{tabularx}
\caption{Métriques du protocole.}
\label{tab:metriques}
\end{table}

\paragraph{Énergie.} Conformément au cahier des charges, l'énergie n'est rapportée que si la mesure
\texttt{powermetrics} a effectivement été réalisée. Sinon, latence, tokens et mémoire sont rapportés
comme indicateurs indirects, sans être convertis en joules.

%======================================================================
\section{Analyse statistique}\label{sec:stats}
%======================================================================

\begin{itemize}
\item \textbf{Intervalles de confiance} à 95\,\% par bootstrap (10\,000 rééchantillonnages des
questions).
\item \textbf{Comparaisons appariées} : les différences entre systèmes sont calculées question par
question, puis rééchantillonnées ensemble (bootstrap apparié), car tous les systèmes répondent aux
mêmes questions.
\item \textbf{Résultats par catégorie} systématiquement rapportés en plus de la moyenne globale.
\item \textbf{Courbes qualité--coût} : score global en fonction de la latence moyenne, pour toutes les
stratégies et pour plusieurs réglages de la méthode adaptative.
\end{itemize}

%======================================================================
\section{Critères de succès (pré-enregistrés)}\label{sec:succes}
%======================================================================

\begin{center}
\fcolorbox{teal!60!black}{teal!5}{\parbox{0.92\textwidth}{
\textbf{Critère principal.} Sur le jeu de test, par rapport à la stratégie B (grand modèle + RAG) :
\begin{enumerate}[label=(\roman*)]
\item la latence moyenne par question est réduite d'\textbf{au moins 20\,\%} ;
\item la perte de score global est d'\textbf{au plus 2 points}, en estimation ponctuelle.
\end{enumerate}
La \emph{non-infériorité} n'est revendiquée que si la borne supérieure de l'intervalle de confiance
à 95\,\% de la perte est elle-même inférieure à 2 points.}}
\end{center}

\paragraph{Critères secondaires.}
\begin{itemize}
\item \textbf{Fiabilité (QR2)} : à couverture égale, le risque de la méthode adaptative est inférieur à
celui de B ; à défaut, la méthode atteint une meilleure couverture à risque égal.
\item \textbf{Contrôle (L1)} : la méthode fait mieux que le routeur fondé sur la seule question. Sinon,
les signaux plus coûteux ne sont pas justifiés, et ce constat sera rapporté.
\item \textbf{Robustesse (QR4)} : aucune catégorie ne perd plus de 5 points par rapport à B. Toute
catégorie qui dépasse ce seuil sera signalée explicitement.
\end{itemize}

\paragraph{Résultat négatif.} Si le critère principal n'est pas atteint, le résultat sera rapporté tel
quel, avec l'analyse des catégories et des signaux responsables. Le cahier des charges rappelle qu'un
résultat négatif bien établi reste informatif.

%======================================================================
\section{Pilote}\label{sec:pilote}
%======================================================================

\paragraph{Déroulement.} Le pilote a exécuté les quatre stratégies fixes sur les 180 questions du jeu de
\emph{développement} (le jeu de test n'a pas été lu), en une passe, le 18 octobre 2026. Il a duré
351~s au total, chargement des modèles et calcul des embeddings du corpus compris. Les sorties brutes
sont dans \texttt{results/pilote/sorties.jsonl} et le résumé dans \texttt{results/pilote/resume.json}.

\begin{table}[H]
\centering\footnotesize\setlength{\tabcolsep}{3.5pt}
\begin{tabular}{l c c c c c c c}
\toprule
\textbf{Stratégie} & \textbf{Score} & \textbf{IC 95\,\%} & \textbf{Couv.} & \textbf{Risque} & \textbf{Lat. moy.} & \textbf{P95} & \textbf{Tokens} \\
\midrule
A -- grand, sans RAG & 0,313 & [0,249 ; 0,377] & 0,31 & 0,82 & 0,25 s & 0,43 s & 62 \\
B -- grand + RAG     & \textbf{0,424} & [0,355 ; 0,493] & 0,56 & 0,53 & 0,99 s & 1,55 s & 579 \\
C -- petit, sans RAG & 0,233 & [0,175 ; 0,295] & 0,46 & 0,93 & 0,13 s & 0,22 s & 84 \\
D -- petit + RAG     & 0,298 & [0,235 ; 0,363] & 0,78 & 0,71 & 0,39 s & 0,59 s & 601 \\
\midrule
\textit{Oracle} & \textit{0,623} & -- & -- & -- & \textit{0,30 s} & -- & -- \\
\bottomrule
\end{tabular}
\caption{Résultats du pilote (jeu de développement, 180 questions). Oracle : stratégie la moins coûteuse donnant le meilleur score pour chaque question. Couverture et
risque : voir tableau~\ref{tab:metriques}. Pic de mémoire MLX du processus, les deux modèles étant chargés :
3,4 à 4,1~Go selon la stratégie.}
\label{tab:pilote}
\end{table}

\begin{table}[H]
\centering\small
\begin{tabular}{l c c c c}
\toprule
\textbf{Score global par catégorie} & \textbf{A} & \textbf{B} & \textbf{C} & \textbf{D} \\
\midrule
Simple       & 0,07 & \textbf{0,45} & 0,05 & 0,33 \\
Complexe     & 0,09 & 0,27 & 0,07 & \textbf{0,31} \\
Non couverte & \textbf{0,78} & 0,55 & 0,58 & 0,25 \\
\bottomrule
\end{tabular}
\hspace{1em}
\begin{tabular}{l c c c c}
\toprule
\textbf{Rappel de la recherche} & \textbf{@1} & \textbf{@3} & \textbf{@6} & \textbf{@10} \\
\midrule
Simple       & 0,48 & 0,60 & 0,70 & 0,77 \\
Complexe     & 0,00 & 0,22 & 0,35 & 0,45 \\
\bottomrule
\end{tabular}
\caption{En haut : score global par catégorie (60 questions chacune). En bas : rappel de la
recherche documentaire seule (pour une question complexe, les deux passages justificatifs doivent
être récupérés).}
\label{tab:pilote-cat}
\end{table}

\paragraph{Enseignements.}
\begin{enumerate}
\item \textbf{La chaîne complète fonctionne} : construction du benchmark, recherche, génération,
notation et journalisation. Les sorties sont déterministes.
\item \textbf{L'écart de coût entre stratégies est important} : B est 7,8 fois plus lente que C et
consomme 7 fois plus de tokens. Le critère principal impose une latence moyenne inférieure à 0,79~s,
contre 0,99~s pour B.
\item \textbf{Aucune stratégie fixe n'est la meilleure partout} : B l'emporte sur les questions simples,
D sur les complexes et A sur les non couvertes. Le routeur \emph{oracle} atteint 0,62 pour 0,30~s par
question. C'est une borne théorique, inatteignable en pratique, mais elle montre une marge réelle pour
une politique adaptative.
\item \textbf{L'abstention laissée au modèle est mal calibrée} : le petit modèle avec RAG répond presque
toujours (couverture 0,78, risque 0,71), tandis que le grand modèle sans RAG s'abstient sur 69\,\% des
questions. Cela justifie de confier la décision d'abstention à la politique (QR2), plutôt qu'au seul
prompt.
\item \textbf{La recherche documentaire est le maillon faible pour les questions complexes} : les deux
passages justificatifs ne sont récupérés ensemble que dans 22\,\% des cas avec $k=3$.
\item \textbf{Le coût de calcul est faible} : une passe complète des quatre stratégies sur le jeu de
test est estimée à environ 15 minutes. Trois répétitions pour la latence sont donc réalistes.
\item \textbf{L'énergie suit les mêmes écarts que la latence} (tableau~\ref{tab:energie}) : une question
coûte 22,8~J avec B contre 1,9~J avec C, soit 12 fois plus. Le RAG augmente à la fois la durée et la
puissance moyenne (environ 24~W contre 17~W), car le traitement du long prompt sollicite davantage le GPU.
\end{enumerate}

\begin{table}[H]
\centering\small
\begin{tabular}{l c c c c}
\toprule
\textbf{Stratégie} & \textbf{Durée} & \textbf{Puissance moy.} & \textbf{Énergie nette} & \textbf{J / question} \\
\midrule
A -- grand, sans RAG & 45,5 s  & 17,1 W & 680 J   & 3,8 \\
B -- grand + RAG     & 181,1 s & 24,8 W & 4\,104 J & \textbf{22,8} \\
C -- petit, sans RAG & 22,9 s  & 16,7 W & 333 J   & \textbf{1,9} \\
D -- petit + RAG     & 71,5 s  & 23,8 W & 1\,547 J & 8,6 \\
\bottomrule
\end{tabular}
\caption{Énergie mesurée par \texttt{powermetrics} (puissance CPU + GPU + ANE, échantillonnée toutes
les 200~ms), 18 octobre 2026, jeu de développement (180 questions). Énergie nette = énergie mesurée
moins la consommation au repos (2,2~W) sur la même durée. La fenêtre de chaque stratégie inclut ses
3 requêtes de chauffe, ce qui surestime d'environ 2\,\% l'énergie par question. Script :
\texttt{bench/energie.py} ; résultats : \texttt{results/pilote/energie.json}.}
\label{tab:energie}
\end{table}

\paragraph{Ajustement prévu avant le gel.} Seul le choix de la méthode de recherche reste ouvert. Avant
de créer l'étiquette \texttt{protocole-v1.0}, on comparera \emph{sur le jeu de développement
uniquement} la recherche dense actuelle, une recherche lexicale BM25 et une combinaison des deux.
Le critère retenu est le rappel@3, à coût comparable. La méthode choisie sera ensuite figée pour
toutes les stratégies.

%======================================================================
\section{Gel du protocole}\label{sec:gel}
%======================================================================

Avant l'évaluation finale :
\begin{enumerate}
\item vérifier que \texttt{data/benchmark/MANIFEST.json} (empreintes des fichiers) n'a pas changé ;
\item figer \texttt{configs/protocole\_v1.json} et créer l'étiquette Git \texttt{protocole-v1.0} ;
\item développer et régler la méthode adaptative uniquement sur \texttt{dev.jsonl} ;
\item exécuter \emph{une seule fois} toutes les méthodes sur \texttt{test.jsonl}, et publier les
sorties brutes dans \texttt{results/}.
\end{enumerate}

\appendix
\section{Prompts}\label{ann:prompts}
\textbf{Sans recherche documentaire :}
\begin{small}
\begin{verbatim}
Answer the question with a short phrase only. If you do not know the answer
or it cannot be determined, reply exactly: UNANSWERABLE.

Question: {question}
Answer:
\end{verbatim}
\end{small}
\textbf{Avec recherche documentaire :}
\begin{small}
\begin{verbatim}
Answer the question using only the passages below, with a short phrase only.
If the passages do not contain the answer, reply exactly: UNANSWERABLE.

[1] {titre}: {texte}
[2] ...

Question: {question}
Answer:
\end{verbatim}
\end{small}

\section{Reproduire}\label{ann:reproduire}
\begin{small}
\begin{verbatim}
python -m venv .venv && source .venv/bin/activate
pip install -r requirements.txt
python bench/build_benchmark.py          # télécharge, vérifie, construit
python bench/check_env.py                # relève l'environnement
HF_HUB_OFFLINE=1 python bench/run_pilot.py   # pilote (jeu de développement)
\end{verbatim}
\end{small}

\bibliographystyle{plainnat}
\bibliography{references}

\end{document}
